\input{../article_base.tex}
\title{פתרון מטלה 01 --- תבניות ריבועיות ומספרים P־אדיים, 80507}

\begin{document}
\maketitle
\maketitleprint[red]

\question{}
יהי שדה $\FF$ ומטריצות $A, B \in M_n(\FF)$ חופפות.

\subquestion{}
נראה כי אם $A$ סימטרית אז $B$ סימטרית.
\begin{proof}
	נניח ש־$P \in M_n(\FF)$ מטריצה הפיכה המקיימת $A = P^t B P$.
	נסמן $g : \FF^n \times \FF^n \to \FF$ התבנית בילינארית הסימטרית המקיימת ${[g]}_E = A$, אז הוכחנו בכיתה כי ${[g]}_P = B$.
	בבניית $g$ השתמשנו בהנחה ש־$A$ סימטרית, ומהנתון החדש על $B$ נקבל שגם היא סימטרית.
\end{proof}

\subquestion{}
נראה כי אם $A$ הפיכה אז $B$ הפיכה.
\begin{proof}
	נניח ש־$P \in M_n(\FF)$ מטריצה הפיכה המקיימת $B = P^t A P$.
	$P^t$ הפיכה אף היא משיקולי דרגה, וכפל של הפיכות משמר הפיכות, לכן גם $A P$ וכן $ A^t (A P)$ הפיכות ובפרט $B$ הפיכה.
\end{proof}

\subquestion{}
נוכיח שקיים $0 \ne c \in \FF$ כך ש־$\det B = c^2 \det A$.
\begin{proof}
	%מהסעיף הקודם נסיק ש־$A$ הפיכה אם ורק אם $B$ הפיכה, ולכן אם $A$ לא הפיכה ו־$\det A = 0$ אז גם $\det B = 0$ ונוכל לבחור $c = 1$.
	נסמן שוב $A = P^t B P$, אז,
	\[
		\det A
		= \det(P^t B P)
		= (\det P^t) (\det B) (\det P)
		= c \cdot \det B \cdot c
		= c^2 \det B
	\]
	כאשר $c = \det P$.
\end{proof}

\question{}
נגדיר את המטריצות,
\[
	A = \begin{pmatrix}
		1 & 0 \\
		0 & 2
	\end{pmatrix},
	I_2 \in M_2(\FF)
\]
ונבדוק אם הן חופפות עבור הגדרות $\FF$.

\subquestion{}
נגדיר $\FF = \RR$ ונראה שהן חופפות.
\begin{proof}
	נבחין כי $2 = 1 \cdot \sqrt{2}^2$ ולכן $A$ חופפת ל־$I_2$ עם הבסיס $(1, 0), (0, \frac{1}{\sqrt{2}})$ מטענה על חפיפה לתת־קבוצת כיסוי.
\end{proof}

\subquestion{}
נגדיר $\FF = \QQ$ ונראה שהן לא חופפות.
\begin{proof}
	אילו הן היו חופפות אז היה באופן דומה לסעיף הקודם מתקיים $2 = 1 \cdot q^2$ עבור $q \in \QQ$ כלשהו, כלומר נקבל $\sqrt{2} \in \QQ$ בסתירה.
\end{proof}

\subquestion{}
נגדיר $\FF = \FF_3$ ונראה שהן לא חופפות.
\begin{proof}
	המטריצות חופפות אם קיים $c \in \FF_3^*$ המקיים $2 = 1 \cdot c^2$, אבל $c = 1 \implies c^2 = 1$ וכן $c = 2 \implies c^2 = 1$ ולכן אין שורש ל־$2$ ובהתאם אין חפיפה.
	נבחין כי $|\FF_3^*| = 2$ ולכן $R = \{1, 2\}$ היא קבוצת הנציגים היחידה שקיימת ונקבל שוב שאין חפיפה.
\end{proof}

\subquestion{}
עבור $\FF = \FF_7$ נראה חפיפה.
\begin{proof}
	נבחין כי $c = 3$ הוא איבר שונה מאפס המקיים $c^2 =_{\FF_7} 2$ ולכן $2 = 1 \cdot c^2$ ונקבל ש־$A$ חופפת ל־$I_2$ עם הבסיס $(1, 0), (0, \frac{1}{3})$.
\end{proof}

\question{}
יהי $V$ מרחב וקטורי מעל $\FF$ ו־$g : V \times V \to \FF, q : V \to \FF$ פונקציות כך שמתקיים,
\[
	\forall v, w \in V,\ 
	g(v, w)
	= \frac{1}{2} (q(v + w) - q(v) - q(w))
\]
נראה שאם $g$ תבנית בילינארית סימטרית וכן $q(c v) = c^2 q(v)$ לכל $c \in \FF, v \in V$ אז $q$ תבנית ריבועית.
\begin{proof}
	נראה טענה מעט חזקה יותר, נראה ש־$q$ היא התבנית הריבועית הנוצרת על־ידי $g$, כלומר,
	\[
		\forall v \in V,\ 
		q(v) = g(v, v)
	\]
	נציב,
	\[
		g(v, v)
		= \frac{1}{2}( q(v + v) - q(v) - q(v))
		\overset{(1)}{=}  \frac{1}{2} (4q(v) - q(v) - q(v))
		= q(v)
	\]
	כאשר $(1)$ נובע מהנתון אודות ההומוגניות של $q$.
	קיבלנו שלפי הגדרה $q$ אכן התבנית הריבועית המתאימה ל־$g$.
\end{proof}

\question{}
נמצא דוגמה למרחב וקטורי $V$ מעל $\FF$ ופונקציות $g : V \times V \to \FF, q : V \to \FF$ כך שמתקיים,
\[
	\forall v, w \in V,\ 
	g(v, w)
	= \frac{1}{2} (q(v + w) - q(v) - q(w))
\]
ו־$g$ תבנית בילינארית סימטרית אבל ש־$q$ איננה תבנית ריבועית.
\begin{solution}
	נגדיר $\FF = \FF^3$ וכן $V = \FF^1$, $q(x) = x^3$.
	בהתאם מתקיים,
	\[
		q(v + w)
		= {(v + w)}^3
		= v^3 + 3 (\cdots) + w^3
		= v^3 + w^3
		= q(v) + q(w)
	\]
	ובהתאם $q(v + w) - q(v) - q(w) = 0$, נקבל אם כך ש־$g(v, w) = 0$ התבנית הקבועה מקיימת את הטענה עבור $q$ שהגדרנו.
	בהתאם אם $q$ הייתה תבנית ריבועית אז היינו מקבלים $1 = 1^3 = q(1) = g(1, 1) = 0$ בסתירה.
\end{solution}

\question{}
יהי $V$ מרחב וקטורי מעל $\FF$ ו־$q : V \to \FF$ תבנית ריבועית, ויהיו $v, w \in V$ כך ש־$q(v) \ne 0$. \\
נראה שקיים $c \in \FF$ יחיד כך ש־$w - c v \in {\{ v \}}^{\perp}$ ונמצא נוסחה מפורשת ל־$c$ כתלות ב־$q, v, w$.
\begin{proof}
	נגדיר $g : V \times V \to \FF$ התבנית הבילינארית היחידה המתאימה ל־$q$, ולכן $g(v, v) \ne 0$.
	מתקיים,
	\[
		w - cv \in {\{ v \}}^\perp
		\iff g(w - cv, v) = 0
		\iff g(w, v) - c g(v, v) = 0
		\iff g(w, v) = c q(v)
	\]
	נגדיר אם כך $c = \frac{g(w, v)}{q(v)}$, הגדרה זו תקפה שכן $q(v) \ne 0$.
	משרשרת הגרירות נקבל שגם $w - c v \in {\{ v \}}^\perp$ כפי שרצינו, ומנוסחה מפורשת זו נקבל גם יחידות.
\end{proof}

\question{}
נמצא דוגמה למרחב וקטורי $V$ מעל $\FF$, תבנית ריבועית $q : V \to \FF$ ווקטורים $v, w \in V$ כך ש־$w - c v \notin {\{ v \}}^\perp$ לכל $c \in \FF$.
\begin{solution}
	נגדיר $\FF = \RR, V = \RR^2$ וכן $q(\begin{pmatrix} x_1 \\ x_2 \end{pmatrix}) = x_1^2 - x_2^2$, ראינו בשיעור כי זוהי אכן תבנית ריבועית. \\
	נקבל שגם $g(\begin{pmatrix} x_1 \\ x_2 \end{pmatrix}, \begin{pmatrix} y_1 \\ y_2 \end{pmatrix}) = x_1 y_1 - x_2 y_2$ היא התבנית הבילינארית הסימטרית המתאימה לה.
	נגדיר $v = (1, -1)$ ולכן $q(v) = 0$, ונגדיר את $w = (1, 0)$, נקבל,
	\[
		w - c v
		= (1, 0) - c (1, -1)
		= (1 - c, -c)
	\]
	ובהתאם,
	\[
		g(v, w - c v)
		= g((1, -1), (1 - c, -1))
		= 1 \cdot (1 - c) - {(-1)}^2
		= 1 - c - c
		= 1 - 2 c
	\]
\end{solution}

\end{document}
